\documentclass[11pt,a4paper]{article}

\usepackage{fontspec}
\usepackage{polyglossia}
\setmainlanguage{russian}
\setotherlanguage{english}

\IfFontExistsTF{Times New Roman}
{\setmainfont{Times New Roman}}
{\setmainfont{Libertinus Serif}}

\usepackage[a4paper,top=1.8cm,bottom=1.8cm,left=2.0cm,right=2.0cm]{geometry}
\usepackage{xcolor}
\usepackage{titlesec}
\usepackage{enumitem}
\usepackage{amsmath,amssymb}
\usepackage{hyperref}

\definecolor{hseblue}{HTML}{144F7B}
\definecolor{softblue}{HTML}{EAF4FA}

\hypersetup{colorlinks=true, linkcolor=black, urlcolor=hseblue}

\setlength{\parindent}{0pt}
\setlength{\parskip}{0.5em}
\linespread{1.05}\selectfont
\setlist[itemize]{leftmargin=1.5em, topsep=2pt, itemsep=2pt}

\titleformat{\section}
  {\Large\bfseries\color{hseblue}}
  {\thesection.}
  {0.5em}
  {}

\titleformat{\subsection}
  {\large\bfseries\color{hseblue}}
  {\thesubsection.}
  {0.5em}
  {}

\newcommand{\speechbox}[1]{%
  \noindent\colorbox{softblue}{%
    \parbox{\dimexpr\linewidth-2\fboxsep\relax}{\textbf{Как сказать устно:} #1}%
  }\par\vspace{0.2em}
}

\begin{document}

\begin{center}
{\Large\bfseries Модели оценки опционов\par}
\vspace{0.35em}
{\large Что сказать, если попросят объяснить Кокса--Росса--Рубинштейна, Блэка--Шоулза--Мертона и Black-76\par}
\vspace{0.45em}
Ковалёва Мария Сергеевна \\
НИУ ВШЭ, Факультет экономических наук, Москва, 2026
\end{center}

\vspace{0.5em}
Этот файл написан как устная шпаргалка. Его можно использовать, если на защите попросят рассказать, что это за модели, чем они отличаются и почему в работе основная эмпирическая модель --- CRR, а Black-76 используется как европейский benchmark.

\tableofcontents
\newpage

\section{Короткий ответ на 1 минуту}

\speechbox{В работе рассматриваются три классические модели оценки опционов. Black--Scholes--Merton --- базовая непрерывная модель для европейских опционов на спот-актив. Она показывает, что цену опциона можно получить через безарбитражное хеджирование и риск-нейтральное ожидание. Black-76 --- близкая по логике формула, но адаптированная для опционов на фьючерсы: вместо спот-цены в ней используется фьючерсная цена. Модель Кокса--Росса--Рубинштейна --- дискретная биномиальная модель: она строит дерево возможных будущих цен и двигается назад от экспирации к текущему моменту. Её главное преимущество для моей работы в том, что она умеет учитывать американское исполнение, то есть право исполнить опцион досрочно. Поэтому CRR используется как основная модель для американских опционов на фьючерс MXI, а Black-76 --- как европейский эталон без раннего исполнения.}

\section{Общая идея моделей оценки опционов}

Все три модели решают одну задачу: определить теоретическую цену опциона при заданных параметрах. У опциона есть базовый актив, страйк, срок до экспирации, процентная ставка, волатильность и тип опциона: call или put. Модель берёт эти параметры и возвращает цену, которую можно сравнить с рыночной ценой.

Главная экономическая идея у всех моделей одна: цена опциона должна быть такой, чтобы не возникало арбитража. Если опцион можно идеально или приближённо воспроизвести портфелем из базового актива и безрискового актива, то цена опциона должна совпадать со стоимостью такого воспроизводящего портфеля.

Вторая важная идея --- риск-нейтральное ценообразование. В реальном мире инвесторы требуют премию за риск, но для оценки производного инструмента можно перейти к специальной риск-нейтральной мере. В ней будущий ожидаемый payoff опциона дисконтируется по безрисковой ставке. Это не значит, что инвесторы действительно нейтральны к риску; это технический способ получить безарбитражную цену.

\section{Black--Scholes--Merton}

\subsection{Что это за модель}

\speechbox{Black--Scholes--Merton --- это классическая модель для оценки европейских опционов на актив, который торгуется на спот-рынке. Она предполагает, что цена базового актива следует геометрическому броуновскому движению, волатильность и ставка постоянны, рынок ликвиден, а опцион можно исполнить только в дату экспирации.}

Модель Black--Scholes--Merton является теоретической базой современной опционной математики. В ней цена базового актива $S_t$ описывается стохастическим процессом:

\[
dS_t=\mu S_t\,dt+\sigma S_t\,dW_t,
\]

где $\mu$ --- ожидаемая доходность, $\sigma$ --- волатильность, $W_t$ --- винеровский процесс. Интуитивно это означает, что доходность актива имеет случайную компоненту, а масштаб этой случайности задаётся волатильностью.

\subsection{Как работает формула}

Для европейского call-опциона без дивидендов формула:

\[
C=S_0N(d_1)-Ke^{-rT}N(d_2),
\]

для европейского put-опциона:

\[
P=Ke^{-rT}N(-d_2)-S_0N(-d_1),
\]

где

\[
d_1=\frac{\ln(S_0/K)+(r+0.5\sigma^2)T}{\sigma\sqrt{T}},
\qquad
d_2=d_1-\sigma\sqrt{T}.
\]

Здесь $S_0$ --- текущая цена базового актива, $K$ --- страйк, $T$ --- срок до экспирации, $r$ --- безрисковая ставка, $\sigma$ --- волатильность, $N(\cdot)$ --- функция распределения стандартной нормальной величины.

\subsection{Как объяснить интуицию}

В формуле есть две части. Для call-опциона $S_0N(d_1)$ можно интерпретировать как приведённую к вероятностной логике часть, связанную с получением базового актива. Вторая часть, $Ke^{-rT}N(d_2)$, связана с приведённой стоимостью выплаты страйка. Разница между ними даёт цену права купить актив по страйку.

Главное, что в формуле нет ожидаемой доходности $\mu$. Это часто спрашивают. Ответ такой: ожидаемая доходность исчезает из-за безарбитражного хеджирования. Цена опциона зависит не от того, какую доходность мы субъективно ждём, а от того, как можно воспроизвести риск опциона через базовый актив и безрисковый актив.

\subsection{Ограничения}

Black--Scholes--Merton плохо описывает реальный рынок, если:

\begin{itemize}
    \item волатильность меняется во времени;
    \item есть улыбка подразумеваемой волатильности;
    \item доходности имеют тяжёлые хвосты и скачки;
    \item опцион можно исполнить досрочно;
    \item базовый инструмент --- фьючерс, а не спот-актив.
\end{itemize}

В моей работе эта модель важна как теоретическая основа: она показывает общую логику безарбитражного ценообразования. Но для эмпирической оценки опционов на фьючерс напрямую используется не Black--Scholes--Merton, а Black-76 и CRR.

\section{Black-76}

\subsection{Что это за модель}

\speechbox{Black-76 --- это модификация логики Black--Scholes для опционов на фьючерсы. Главное отличие в том, что в формулу входит не спот-цена актива $S$, а фьючерсная цена $F$. Модель остаётся европейской: она оценивает опцион, который исполняется только на экспирации.}

Black-76 особенно удобна для фьючерсных опционов, потому что в них базовым активом является фьючерсный контракт. В моей работе опционы написаны на фьючерс MXI, поэтому если использовать аналитическую европейскую формулу, корректнее брать именно Black-76.

\subsection{Формулы}

Для call-опциона на фьючерс:

\[
C=e^{-rT}\left(FN(d_1)-KN(d_2)\right),
\]

для put-опциона:

\[
P=e^{-rT}\left(KN(-d_2)-FN(-d_1)\right),
\]

где

\[
d_1=\frac{\ln(F/K)+0.5\sigma^2T}{\sigma\sqrt{T}},
\qquad
d_2=d_1-\sigma\sqrt{T}.
\]

Здесь $F$ --- текущая фьючерсная цена. В отличие от Black--Scholes--Merton, в числителе $d_1$ нет члена $rT$ рядом с логарифмом, потому что фьючерсная цена уже отражает стоимость переноса базового актива.

\subsection{Как объяснить отличие от Black--Scholes--Merton}

Самый простой ответ: Black--Scholes--Merton оценивает опцион на спот-актив, а Black-76 --- опцион на фьючерс. Если базовый актив --- фьючерс, то использовать $F$ логичнее, чем $S$, потому что именно фьючерс является поставляемым инструментом по опциону.

Экономически Black-76 сохраняет ту же безарбитражную идею: цена равна дисконтированному риск-нейтральному ожиданию payoff. Но динамика строится вокруг фьючерсной цены.

\subsection{Почему Black-76 в работе не основная модель}

Black-76 --- европейская модель. Она не учитывает досрочное исполнение. А рассматриваемые опционы допускают американское исполнение, то есть держатель может исполнить их до экспирации.

Поэтому в моей работе Black-76 используется как benchmark: она показывает, что будет, если игнорировать право раннего исполнения. Результат оказался важным: Black-76 даёт существенно большую ошибку, чем американская CRR-модель. Это подтверждает, что для этих контрактов стиль исполнения является содержательным фактором цены.

\section{Модель Кокса--Росса--Рубинштейна}

\subsection{Что это за модель}

\speechbox{Модель Кокса--Росса--Рубинштейна, или CRR, --- это биномиальная модель. Она разбивает время до экспирации на шаги и предполагает, что на каждом шаге цена базового актива может либо вырасти, либо упасть. Затем модель считает payoff в конце дерева и обратной индукцией возвращается к текущей цене опциона.}

CRR является дискретной моделью. Вместо непрерывного процесса, как в Black--Scholes--Merton, она строит дерево возможных цен. При большом числе шагов биномиальная логика приближается к непрерывной модели, но при этом CRR гибче в одном важном аспекте: в ней удобно учитывать досрочное исполнение.

\subsection{Как строится дерево}

Время до экспирации делится на $N$ шагов:

\[
\Delta t=\frac{T}{N}.
\]

На каждом шаге цена может вырасти в $u$ раз или упасть в $d$ раз:

\[
u=e^{\sigma\sqrt{\Delta t}}, \qquad d=e^{-\sigma\sqrt{\Delta t}}=\frac{1}{u}.
\]

В классической модели для акции риск-нейтральная вероятность:

\[
q=\frac{e^{r\Delta t}-d}{u-d}.
\]

Для фьючерса в моей постановке используется версия:

\[
q=\frac{1-d}{u-d},
\]

потому что в риск-нейтральной мере фьючерсная цена является мартингалом: её ожидаемое значение на следующем шаге равно текущему значению.

\subsection{Как считается цена опциона}

Сначала в конечных узлах дерева считается payoff. Для call:

\[
C_T=\max(F_T-K,0),
\]

для put:

\[
P_T=\max(K-F_T,0).
\]

Затем модель идёт назад по дереву. Для европейского опциона стоимость в узле равна дисконтированному ожидаемому значению будущих стоимостей:

\[
V=e^{-r\Delta t}\left(qV_u+(1-q)V_d\right).
\]

Для американского опциона добавляется выбор: продолжить держать опцион или исполнить его прямо сейчас:

\[
V=\max\left(H(F,K),\ e^{-r\Delta t}\left(qV_u+(1-q)V_d\right)\right),
\]

где $H(F,K)$ --- внутренняя стоимость опциона. Именно этот максимум делает модель подходящей для американских опционов.

\subsection{Как объяснить обратную индукцию}

\speechbox{Обратная индукция означает, что мы сначала знаем стоимость опциона в конце, потому что там она равна payoff. Потом шаг за шагом возвращаемся назад: в каждом узле спрашиваем, сколько стоит право попасть в два будущих состояния, и для американского опциона дополнительно сравниваем это со стоимостью немедленного исполнения.}

Это удобно объяснять так: в конце дерева неопределённости уже нет, поэтому цена опциона очевидна. До экспирации цена равна стоимости выбора между будущими сценариями. Для американского опциона выбор шире, потому что можно не ждать будущего, а исполнить опцион сейчас.

\subsection{Почему CRR основная модель в работе}

В моей работе исследуются американские опционы на фьючерс MXI. Значит, модель должна учитывать:

\begin{itemize}
    \item базовый актив как фьючерсную цену;
    \item разные способы задания волатильности;
    \item возможность досрочного исполнения.
\end{itemize}

CRR закрывает эти требования. Поэтому я использую три CRR-спецификации: с HV\_21d, HV\_63d и GARCH-волатильностью. Black-76 остаётся европейским benchmark, чтобы отдельно оценить стоимость раннего исполнения.

\section{Как сравнить три модели между собой}

\begin{itemize}
    \item \textbf{Black--Scholes--Merton} --- базовая непрерывная модель для европейских опционов на спот-актив.
    \item \textbf{Black-76} --- европейская формула для опционов на фьючерсы.
    \item \textbf{CRR} --- биномиальная модель, которая подходит и для европейских, и для американских опционов, потому что допускает проверку досрочного исполнения в каждом узле.
\end{itemize}

\speechbox{Если совсем кратко: Black--Scholes--Merton даёт теоретическую основу, Black-76 адаптирует эту основу к фьючерсам, а CRR позволяет учесть американский стиль исполнения. Поэтому для моей эмпирической задачи CRR является основной моделью, а Black-76 нужен как контрольный европейский ориентир.}

\section{Готовые ответы на вопросы}

\subsection*{Почему нельзя было использовать только Black--Scholes--Merton?}

Потому что Black--Scholes--Merton предназначена для европейских опционов на спот-актив. В моей работе опционы написаны на фьючерс и имеют американский стиль исполнения. Поэтому напрямую BSM не соответствует инструменту. Корректнее использовать Black-76 для фьючерсной европейской постановки и CRR для американской постановки.

\subsection*{Почему Black-76 хуже CRR в результатах?}

Не потому что формула «плохая», а потому что она решает другую задачу. Black-76 оценивает европейский опцион, а исследуемые контракты допускают досрочное исполнение. CRR учитывает это право, поэтому для американских опционов она ближе к экономической структуре инструмента.

\subsection*{Что значит «досрочное исполнение»?}

Это право держателя американского опциона исполнить контракт до даты экспирации. В CRR на каждом узле дерева модель сравнивает стоимость продолжения со стоимостью немедленного исполнения и выбирает максимум.

\subsection*{Почему волатильность так важна во всех моделях?}

Опцион ценен из-за неопределённости. Чем выше волатильность, тем больше вероятность, что payoff окажется выгодным для держателя. Поэтому разные оценки волатильности дают разные модельные цены.

\subsection*{Почему в CRR можно использовать разные волатильности?}

CRR сама задаёт механизм ценообразования, но ей нужен входной параметр $\sigma$. В моей работе этот параметр задаётся тремя способами: через 21-дневную историческую волатильность, 63-дневную историческую волатильность и GARCH. Поэтому сравниваются не только модели цены, но и связки «модель цены плюс модель волатильности».

\subsection*{Какой главный вывод по моделям?}

Главный вывод: сначала нужно правильно выбрать класс модели под инструмент. Для американского опциона на фьючерс важнее учесть раннее исполнение, чем просто подобрать более сложную волатильность. Поэтому CRR методологически является основной моделью в работе.

\section{Финальная формулировка для защиты}

Если комиссия попросит рассказать про модели целиком, можно ответить так:

\textit{В работе я опираюсь на три классические модели. Black--Scholes--Merton --- это базовая непрерывная модель для европейских опционов на спот-актив; она важна как теоретическая основа безарбитражного ценообразования. Black-76 сохраняет похожую логику, но применяется к опционам на фьючерсы, поэтому вместо спот-цены использует фьючерсную цену. Однако обе эти аналитические модели являются европейскими и не учитывают досрочное исполнение. Для моих данных это критично, потому что опционы на фьючерс MXI имеют американский стиль исполнения. Поэтому основная модель в эмпирической части --- биномиальная модель Кокса--Росса--Рубинштейна: она строит дерево возможных цен, считает payoff на экспирации и обратной индукцией возвращается к текущей цене, на каждом шаге сравнивая продолжение с немедленным исполнением. Именно поэтому CRR подходит для основной оценки, а Black-76 используется как benchmark, показывающий, насколько много мы теряем, если игнорируем право раннего исполнения.}

\end{document}
